\subsection{A framed combinatorial criterion for the arithmetic image}

We now remove the conjugacy ambiguity in E.4 and combine it with E.13.
The result addresses cote 142, manuscript p.~8, scan~17, together with
the covering and puncture-filling structures on pp.~8v--10, scans~18--21.
We fix the standard genus-zero tangential marking: the two positively
oriented generators $x,y$ of the profinite free group
$\widehat F_2$ are the loops about $0,1$ based at the tangent from $0$
towards $1$, with the real interval fixing the path convention. In the
combinatorial presentation this is simply a choice of generators and
of the reflection
\[
             \iota(x)=x^{-1},\qquad \iota(y)=y^{-1}.
\]
The algebraic interpretation verifies the criterion; it is not needed
to state its tests.

Write $B=\widehat{\mathrm{GT}}$ in the convention of
\citet[Definition, p.~572; Theorem~1, p.~573]{E_LochakSchneps1997}.
Concretely its elements are the invertible pairs $(\lambda,f)$, with
$\lambda\in\widehat\Z^{\times}$ and
$f\in\overline{[\widehat F_2,\widehat F_2]}$, satisfying the symmetry,
hexagon and pentagon relations. Here invertible means that
\[
       x\longmapsto x^\lambda,
       \qquad y\longmapsto f^{-1}y^\lambda f
                                                               \tag{E.14.1}
\]
is an automorphism, and the relations, with $z=(xy)^{-1}$ and
$m=(\lambda-1)/2$, are
\[
\begin{split}
 f(y,x)f(x,y)&=1,\\
 f(z,x)z^m f(y,z)y^m f(x,y)x^m&=1,\\
 f(x_{34},x_{45})f(x_{51},x_{12})f(x_{23},x_{34})
 f(x_{45},x_{51})f(x_{12},x_{23})&=1.
\end{split}                                                     \tag{E.14.2}
\]
The last equation is in the pure five-punctured-sphere group,
presented by its usual pure braid generators $x_{ij}$.
The product is
\[
 (\lambda,f)(\mu,g)=
 \bigl(\lambda\mu,\,f\,g(x^\lambda,f^{-1}y^\lambda f)\bigr).
                                                               \tag{E.14.3}
\]
In particular $c=(-1,1)$ is a specified nonidentity involution of $B$.

There is a completely specified genus-zero outer-action system
\[
 \tau_0:B\longrightarrow
 \prod_{n,\,\mathcal P}\operatorname{Out}
                    (\widehat M_{0,n,\mathcal P}),             \tag{E.14.4}
\]
where $\mathcal P$ records any partition into labels and colours.
For clarity, it can be computed on the Artin braid groups as follows.
Put $y_1=1$ and
$y_i=\sigma_{i-1}\cdots\sigma_2\sigma_1^2
\sigma_2\cdots\sigma_{i-1}$. The action of $(\lambda,f)$ is
\[
 \sigma_1\longmapsto\sigma_1^\lambda,
 \qquad
 \sigma_i\longmapsto
 f(y_i,\sigma_i^2)^{-1}\sigma_i^\lambda f(y_i,\sigma_i^2).
                                                               \tag{E.14.5}
\]
Pass to the sphere quotient $y_n=\omega_n=1$, with
$\omega_n=y_1\cdots y_n$, and restrict to the inverse images of the
colour-preserving permutation subgroups. This passage really is defined:
the action sends $y_i$ to $y_i^\lambda$ and $\omega_n$ to
$\omega_n^\lambda$, so it and its inverse preserve the closed normal
subgroup generated by the two sphere relations
\citep[Appendix, author p.~35]{E_HarbaterSchneps2000}.
It preserves the permutation quotient, since $\lambda$ is odd and
the conjugating words are pure. Thus it preserves every indicated
colour subgroup. Each pure meridian $x_{ij}$ is a braid conjugate of
some $\sigma_k^2$; preservation of permutations makes its image a
\emph{pure}-braid conjugate of $x_{ij}^{\lambda}$.
Deleting the last distinguished strand also commutes
with the formula: in its permutation stabilizer the deletion kernel is
normally generated by the pure meridians $x_{in}$; these are preserved
up to conjugacy and unit power, while the remaining generators
$\sigma_1,\ldots,\sigma_{n-2}$ have exactly the lower-rank formula.
The sphere quotients inherit this compatibility. Relabelling gives
the other deletions up to the inner conjugacies appropriate for outer
actions. These are the full
profinite braid actions of \citet[\S2, Main Theorem~(iii), equation~(5);
\S5, Propositions~10--11]{E_LochakSchneps1994}; no pro-$\ell$
substitution is made. We use their action, not a claim that every
abstract genus-zero outer system is in its image.

Let $\mathfrak M_{\mathrm{AA}}$ be the set of maximal closed
subgroups of $B$ satisfying exactly the AA-property in
\citet[Definition~3.3(v); Theorem~4.4(ii); Definition~5.12]{E_HMT2025}.
AA means \emph{algebraically ample}. The adjective
``arithmetically ample'' in the earlier audit was an incorrect
expansion of this abbreviation. We retain all the prerequisites:
COF and RGC for the arithmetic Belyi diagrams, ArBC for the
group-constructed operations, and all four clauses of Definition~5.12.
In particular the function-field/Kummer conditions in its clause~(iii)
and the second-configuration compatibility in clause~(iv) are included.
The input second-configuration group is the full profinite group of
the five-punctured-sphere moduli space, with the indicated forgetful
quotient. It is part of the surface presentation already supplied in E.8.
These tests concern objects constructed from profinite groups; they
do not use a prior embedding of an absolute Galois group as a test.

\begin{claim}{E.14}{Proved}
There is exactly one member $G_c$ of $\mathfrak M_{\mathrm{AA}}$
that contains $c=(-1,1)$.

Let $a=(a_{g,n})$ be a family of outer continuous automorphisms of
the full profinite mapping-class groups, in the fixed markings.
It comes from an element of $G_\Q$ if and only if both of the
following purely combinatorial conditions hold:
\begin{enumerate}[label=(\roman*)]
\item The family admits a coherent extension to a continuous,
type-preserving autoequivalence $F$ of the labelled/coloured completed
cover-and-filling presentation of E.8. The extension preserves the
covering and filling generators, the finite permutation quotients,
the label-forgetting maps and all specified puncture-deletion maps.
\item Such an extension can be chosen with its \emph{complete}
genus-zero outer-action system equal to $\tau_0(b)$ for some $b\in G_c$.
\end{enumerate}
For a system already supplied with this coherent structure, condition
(i) is part of its data and only (ii) remains to be tested. The element
of $G_\Q$ producing an arithmetic system is unique.
\end{claim}
\begin{proof}
First we isolate the elementary selection argument. Let $G\subset B$
denote the standard arithmetic image for the duration of the proof.
By \citet[Theorem~1]{E_LochakSchneps1997}, $c$ belongs to $G$ and
is the image of complex conjugation. Their Proposition~4(ii), proved
on printed p.~581, states
\[
                         C_B(c)=\langle c\rangle.             \tag{E.14.6}
\]
Its exact scope is the full profinite group. In their proof, commuting
with $c$ in (E.14.3) forces $f(x,y)=f(x^{-1},y^{-1})$;
Proposition~4(i) makes $f=1$, and the hexagon forces $\lambda=\pm1$.

All nonidentity involutions of $G_\Q$ are conjugate in $G_\Q$.
Indeed their fixed fields in $\overline\Q$ are real closures by the
Artin--Schreier theorem; $\Q$ has one ordering, and its real closures
are isomorphic over $\Q$. Such an isomorphism extends to the algebraic
closures and conjugates the two order-two stabilizers. This is also
the precise involution fact used in the proof of
\citet[Proposition~4(iii), p.~581]{E_LochakSchneps1997}.

Now let $H=bGb^{-1}$ and suppose $c\in H$. Then $b^{-1}cb\in G$
is a nonidentity involution, so $b^{-1}cb=g c g^{-1}$ for some
$g\in G$. Consequently $bg\in C_B(c)=\langle c\rangle\subset G$.
Thus $b\in G$ and $H=G$. By E.4, or directly
\citet[Theorem~5.17(ii)]{E_HMT2025}, the members of
$\mathfrak M_{\mathrm{AA}}$ are exactly these conjugates of $G$.
The subgroup $G$ itself satisfies AA by their Remark~5.12.1 and is
maximal by their Lemma~5.16. Existence and uniqueness follow, and
\[
                              G_c=G.                          \tag{E.14.7}
\]
This equality proves that the group-theoretic selection actually picks
the specified arithmetic subgroup. It is not used in the definition
of the selection procedure.

For necessity, the algebraic presentation of E.8 is defined over
$\Q$. An element $\sigma\in G_\Q$ therefore acts continuously
on its objects, paths, finite covers and puncture fillings, and
preserves their labelled and coloured versions. It induces the
coherent extension in (i). We check the complete genus-zero
comparison, including its finite colour extensions.

Write $K_n=\widehat M_{0,n}^{\rm pure}$. In
\citet[\S0.1, Definition; \S3.1, Proposition~8, Lemma~1 and
Theorem~5, author pp.~2, 29--31]{E_HarbaterSchneps2000},
$\operatorname{Out}^{\sharp}_n$ consists of the outer automorphisms
of $K_n$ that send each pure collision meridian to a conjugate of
its $\lambda$-power, for one $\lambda\in\widehat\Z^\times$, and
commute with the outer $S_n$-action. The standard action (E.14.5)
lands in these groups and respects deletion. Every deletion
$\operatorname{Out}^{\sharp}_n\to
\operatorname{Out}^{\sharp}_{n-1}$ is injective for $n\ge5$.
The arithmetic action also lies there: the collision boundary divisors
are defined over $\Q$, so inertia transforms by the cyclotomic
character, and all label permutations are defined over $\Q$.
The arithmetic and standard actions agree on $K_4$ by the normalized
Ihara--Belyi embedding used in (E.14.1). Their compatibility with
deletion and the cited injectivity give equality on every $K_n$.

For $n\ge5$, the centralizer of $K_n$ inside the full group
$\widehat M_{0,n}$ is trivial. Indeed its permutation acts trivially
on $H_1(K_n,\Q_\ell)$, so is the identity by E.13's faithful
permutation calculation, and the remaining element lies in the
centre-free group $K_n$. Choose the arithmetic representatives using
an ordered rational configuration, so they preserve every colour
subgroup; changing to the standard marking changes them by inner
automorphisms. Since the arithmetic and standard restrictions have the
same outer action on $K_n$, a conjugation by an element of $K_n$
makes them agree there. If their quotient is $\alpha$, then for any
$h$ in the full group, $h^{-1}\alpha(h)$ centralizes $K_n$.
Consequently $\alpha(h)=h$. The same conjugation belongs to every
colour subgroup, proving equality of the outer actions on all of them.
For $n=3,4$, add distinguished punctures to reach at least five and
then use the surjective deletion homomorphism, with its already
verified compatibility. This proves the complete comparison with
$\tau_0(\sigma)$. Since the GT image of $\sigma$ belongs to
$G=G_c$, condition (ii) follows.

There is also a direct check of the framed all-$n$ arithmetic formula
in \citet[\S5.4, equation~(5.4.2), printed p.~28]{E_Nakamura2010},
using the Ihara--Matsumoto tangential configuration
$(0,t^{n-1},\ldots,t)$. That formula uses $\omega_i$ in place of
$y_i$. The identity $\omega_i=y_i\omega_{i-1}$ gives the same
formula: $\omega_{i-1}$ commutes with both $y_i$ and $\sigma_i$,
so the commutator word $f$ ignores that central factor in its first
argument. This corroborates the full-group comparison above; no
pro-$\ell$ or meta-abelian replacement is involved.

Conversely, suppose (i)--(ii) hold. By (E.14.7), the element
$b\in G_c$ is the image of some $\sigma\in G_\Q$. Its arithmetic
action supplies a structured autoequivalence $F_\sigma$ on all
types. The composite $F_\sigma^{-1}F$ has the identity outer action
on every genus-zero group, including every labelled and coloured
version: equality in (ii) is equality of the whole system, not only
of its value on one free group. E.13 consequently gives the identity
outer action on every $\widehat M_{g,n}$. Thus $a$ is exactly the
arithmetic outer-action family of $\sigma$.

Uniqueness also holds when the supplied family records only the full
unlabelled mapping-class groups. Suppose that an arithmetic element
acts inner on $\widehat M_{0,5}$. On its pure subgroup $K_5$ its
outer action is therefore the action of some permutation
$\pi\in S_5$. Arithmetic outer actions commute with $S_5$, and
$S_5\to\operatorname{Out}(K_5)$ is injective by E.13's faithful
first-homology calculation. Thus $\pi$ is central in $S_5$, so
$\pi=1$. The pure outer action is trivial, and deleting one labelled
point gives trivial arithmetic outer action on
$K_4=\widehat F_2$.

This last arithmetic action is faithful. To recall
the Belyi argument, an element acting inner fixes the isomorphism class
of every connected finite cover of the tripod and hence of its smooth
compactification. Every elliptic curve over $\overline\Q$ has such
a Belyi presentation, so all its $j$-invariants are fixed. Every
element of $\overline\Q$ occurs as such a $j$-invariant. The Galois
element is therefore the identity. This proves the asserted uniqueness.
\end{proof}

The criterion has explicit marking dependence. Replacing the standard
genus-zero framing by a $B$-conjugate transports $c$, the set of
maximal AA subgroups, and the selected subgroup together. If the
change preserves $c$, uniqueness preserves $G_c$; in particular a
$B$-inner change fixing $c$ is governed by (E.14.6). On the remaining
objects, changing markings gives the usual inner conjugacy ambiguity
already removed by taking outer automorphism classes. No unchosen
identification of a subgroup with its conjugacy class is made.

Condition (i) is an actual condition on the surface presentation:
choose continuous bijections of its Hom spaces respecting sources,
targets, composition, the specified finite-cover groupoids and the
label/deletion structure, and require their induced object-group
automorphisms to represent the given outer family. All these Hom
spaces and relations were constructed in E.5--E.8. In particular this
condition does not quantify over algebraic curves or over Galois
elements. We claim a mathematical characterization, not a terminating
algorithm for deciding membership in these infinite profinite groups.
The proof neither asserts $B=G_\Q$ nor lifts an arbitrary $B$
conjugator or an arbitrary genus-zero system to all types.
